\begin{officialform}[BulletinVI]
\FormHeading
\FormLine{Mã số đề tài: \ProjectCode}
\FormLine{Tên đề tài: \ProjectTitle}
\FormLine{Thời gian thực hiện: \ProjectPeriod}
\FormLine{Tổng kinh phí: \ProjectBudget}
\FormLine{Chủ nhiệm đề tài: \PrincipalName\ (\PrincipalPhone)}
\FormLine{Thành viên tham gia nghiên cứu: Nguyễn Thanh Trọng, Cao Tường Hưng, Nguyễn Như Quỳnh, Ngô Hưng Thịnh.}
\FormLine{Cán bộ hướng dẫn: \SupervisorName\ (\SupervisorEmail)}

\FormLine{Tính cấp thiết: Trong thiên tai bão lũ, thông tin cầu cứu thường phân tán, trùng lặp và thiếu dữ liệu cần thiết, gây khó khăn cho việc xác minh và điều phối. Hạ tầng viễn thông có thể quá tải hoặc gián đoạn, hạn chế khả năng gửi hình ảnh và tiếp nhận yêu cầu kịp thời. Vì vậy, cần nghiên cứu hệ thống phân tích thông tin ngay trên thiết bị, truyền dữ liệu tối giản và tổng hợp các báo cáo liên quan để hỗ trợ lực lượng cứu hộ nắm bắt tình hình.}

\FormLine{Mục tiêu: Xây dựng hệ thống hỗ trợ tiếp nhận, phân tích và tổng hợp thông tin cứu hộ bão lũ trong điều kiện mạng không ổn định. Kết hợp xử lý AI trên thiết bị di động, truyền dữ liệu tối giản và phân cụm sự kiện nhằm hỗ trợ xác minh, xác định mức độ ưu tiên và trực quan hóa yêu cầu cứu hộ trên bản đồ.}

\FormLine{Phương pháp nghiên cứu: Kết hợp nghiên cứu tài liệu, xây dựng và chuẩn hóa dữ liệu, huấn luyện và tối ưu mô hình AI cho thiết bị di động. Thiết kế ứng dụng, máy chủ và giao diện quản lý; thực nghiệm đánh giá chất lượng nhận diện, phân cụm, xếp hạng ưu tiên, tốc độ xử lý và khả năng truyền dữ liệu trong điều kiện mạng yếu giả lập.}

\FormLine{Nội dung nghiên cứu: Xây dựng và tối ưu mô hình nhận diện ảnh cứu hộ trên thiết bị di động; thu thập thông tin về vị trí và mức độ khẩn cấp; phát triển cơ chế lưu trữ, gửi lại và đồng bộ dữ liệu khi kết nối mạng không ổn định. Nghiên cứu phân cụm, ưu tiên và bản đồ. Bổ sung thiết kế DistilBERT/fusion nhận diện yêu cầu hỗ trợ trong tin tiếng Việt, ablation và đánh giá camera/Android; chưa triển khai hoặc đo các phương án này.}

\FormLine{Kết quả đạt được: Mô hình gốc seed 1024 đạt Accuracy \NewImageAccuracy\%, macro-F1 \NewImageMacroFOne\% trên 256 ảnh CPU, có 3 lỗi nặng. Lượng tử hóa giảm dung lượng và độ chính xác; có 32 phiên benchmark ảnh trên bốn điện thoại. LR năm cờ không phải NLP; adapter chưa dùng điểm LR. Đã có ứng dụng, backend và dashboard. Pilot CrisisMMD tiếng Anh đạt macro-F1 fusion \MMFusionFOne\%, nhưng CI hiệu còn chứa 0; chưa xác nhận NLP tiếng Việt hoặc fusion trong ứng dụng.}

\FormLine{Ý nghĩa: RQ1 chưa xác lập Product hơn Additive; RQ2 chưa xác lập lợi ích alignment chung; RQ3 cho revised kém nearest-first về harm và trễ hạn mô phỏng. Ba raw p Spearman và năm Holm RQ2 chưa khớp replay. Chưa đo hiệu quả cứu hộ thực địa.}

\FormLine{Khả năng ứng dụng: Sử dụng trong giảng dạy, nghiên cứu và trình diễn quy trình quản lý thông tin cứu hộ. Hệ thống là cơ sở để tiếp tục hoàn thiện, phối hợp với đơn vị chuyên môn thử nghiệm và đánh giá trước khi áp dụng thực tế.}
\end{officialform}
